\begin{abstract}
While scaling laws for large language models (LLMs) during pre-training have been extensively studied, their behavior under reinforcement learning (RL) post-training remains largely unexplored. This paper investigates the scaling behavior of Large Language Model (LLM) reinforcement learning post-training, focusing on mathematical reasoning. Through experiments across the Qwen2.5 series (0.5B to 72B), we characterize how model scale, data, and compute interact. Our analysis yields four key findings: \ding{182} Larger models consistently demonstrate superior compute and data efficiency. \ding{183} The relationship between model performance and training resources follows a \textbf{predictive power-law} across both base and instruction-tuned models. \ding{184} RL learning efficiency exhibits a latent \textbf{saturation trend} with increasing model scale.
\ding {185} In data-constrained regimes, performance is primarily driven by the \textbf{total volume of training data} rather than sample uniqueness. These results offer practical guidelines for scaling reasoning capabilities through reinforcement learning post-training.
\end{abstract}